\section{Limits from rules}
\label{sec:limits}

The characterization measures how fast the boat \emph{can} go. That is no speed to plan gates
at: what limits a lagging boat is how far it travels before a new command takes effect. Each
limit the characterization writes to the Nav2 overlay is therefore
\[
  \text{limit} = \min(\text{measured} \times \text{headroom},\; \text{rule}(T)),
\]
with the rules below, in terms of the lag $T_a = \tau_a + L$ of Section~\ref{sec:latency}
(Figure~\ref{fig:limits}).

\subsection{Acceleration: a ramp lasts $k$ lags}
A first-order lag $T$ following a velocity ramp of slope $a$ trails it by $aT$ in steady state.
Ramping to $v$ at $a = v/(kT)$ keeps that tracking error at $v/k$, a $k$-th of the speed, so
MPPI's kinematic plan stays one the boat can follow:
\[
  a_a = \frac{v_a}{k\,T_a},\qquad k = 3 \;(\code{ramp\_lags}).
\]
The same rule sets the deceleration and, with $w_z$, the angular acceleration
$\alpha_z = w_z/(k T_\text{yaw})$.

\subsection{Speed: stop inside $D$}
After a stop command the boat keeps $v$ for about one lag, then decelerates at $a$:
\[
  s_\text{stop} = vT + \frac{v^2}{2a} = vT + \frac{v k T}{2} = vT\Big(\frac{k}{2} + 1\Big) \le D
  \quad\Longrightarrow\quad
  v_a = \frac{D}{T_a\,(k/2 + 1)},
\]
with $D = 1$\,m (\code{stop\_distance\_m}), so an obstacle or a goal one metre ahead is always
stoppable for.

\subsection{Yaw rate: overshoot within half the tolerance}
A turn at rate $w_z$ keeps turning for about one lag after the command to stop, so it
overshoots by $\theta_\text{os} \approx w_z T_\text{yaw}$. Keeping that within half the precise
goal checker's yaw tolerance leaves the other half to heading noise, the same factor of two as
every threshold against the noise it guards:
\[
  w_z\,T_\text{yaw} \le \tfrac12\,\psi_\text{tol}
  \quad\Longrightarrow\quad
  w_z = \frac{\psi_\text{tol}}{2\,T_\text{yaw}},\qquad \psi_\text{tol} = 0.25\text{\,rad}.
\]
The characterization writes it to \code{wz\_max} and the behaviour server's
\code{max\_rotational\_vel}; the rotation shim's \code{rotate\_to\_heading\_angular\_vel} is set
to match by hand. At $L = 0.3$\,s ($T_\text{yaw} = 0.45$\,s) a turn at 0.6\,rad/s would overshoot
about $0.27$\,rad (15$^\circ$), twice the half-tolerance of 7$^\circ$.

\begin{table}[h]
\centering\small
\begin{tabular}{@{}lcccc@{}}
\toprule
axis & $T$ (s) & rule $v$ or $w_z$ & rule $a$ or $\alpha_z$ & today's $v$ or $w_z$ \\
\midrule
surge & $0.35 + 0.15$ & 0.80\,m/s & 0.53\,m/s$^2$ & 0.4\,m/s \\
sway  & $0.45 + 0.15$ & 0.67\,m/s & 0.37\,m/s$^2$ & 0.15\,m/s \\
yaw   & $0.15 + 0.15$ & 0.42\,rad/s & 0.46\,rad/s$^2$ & 0.3\,rad/s \\
\bottomrule
\end{tabular}
\caption{Rule outputs at the \code{nav2\_rules} constants ($\tau$ per axis $+$ \code{delay\_s} 0.15,
$D = 1$\,m, $k = 3$, $\psi_\text{tol} = 0.25$), against today's limits. In the Python plant the caps
cut route time 82.5$\to$47.6\,s calm and 84.4$\to$48.3\,s in the field mix, with goal overshoot
no worse (0.58$\to$0.50\,m in the field mix) \ev{rules\_v1:cur\_calm, rule\_calm, cur\_field, rule\_field}. Under MPPI the yaw caps alone (A3) are not yet accepted:
path ratio 1.32 and one away-spin \ev{tune\_A3\_pair}; the full set awaits the crane matrix.}
\label{tab:rules}
\end{table}

\noindent MPPI's \code{prune\_distance} must cover the horizon at the new speed:
$\ge 1.2\,v_x \times 4.8$\,s. A current $c_\text{max}$ needs $v_x \ge c_\text{max} + 0.2$\,m/s of
headroom, or the day is a no-go.

\begin{figure}[tbp]
\centering
\begin{tikzpicture}
\begin{groupplot}[group style={group size=2 by 1, horizontal sep=1.6cm},
  height=5.2cm, tick label style={font=\footnotesize}, label style={font=\small},
  legend cell align=left,
  legend style={font=\footnotesize, draw=none, fill=white, fill opacity=0.85, text opacity=1}]
% (a) stopping geometry, in units of the lag T and the speed v (k = 3)
\nextgroupplot[width=0.43\linewidth, xmin=-0.3, xmax=4.6, ymin=0, ymax=1.45,
  axis lines=left, xtick={0,1,4}, xticklabels={$0$, $T$, $(k+1)T$}, ytick={1},
  yticklabels={$v$}, xlabel={time after the stop command}, ylabel={speed},
  title={\small (a) stopping takes $vT + v^2/2a$}, clip=false]
  \fill[sA!15] (axis cs:0,0) -- (axis cs:0,1) -- (axis cs:1,1) -- (axis cs:4,0) -- cycle;
  \draw[sA!50, dashed] (axis cs:1,0) -- (axis cs:1,1);
  \addplot[sA, very thick] coordinates {(-0.3,1) (1,1) (4,0) (4.6,0)};
  \node[font=\footnotesize] at (axis cs:0.5,0.5) {$vT$};
  \node[font=\footnotesize] at (axis cs:1.9,0.33) {$\dfrac{v^2}{2a}$};
  \node[font=\scriptsize, anchor=south west, align=left] at (axis cs:2.45,0.55)
    {ramp $a = v/(kT)$};
  \node[font=\footnotesize, anchor=south, align=center] at (axis cs:2.0,1.05)
    {$vT\,(k/2 + 1) \le D$};
% (b) the rule caps against the lag T
\nextgroupplot[width=0.5\linewidth, xmin=0.2, xmax=1.0, ymin=0, ymax=1.2, clip mode=individual,
  xlabel={lag $T = \tau + L$ (s)}, ylabel={cap (m/s, rad/s)},
  title={\small (b) caps from $D = 1$\,m, $k = 3$, tol $= 0.25$\,rad},
  grid=major, grid style={gray!20}, legend pos=north east]
  \addplot[sA, very thick] table[x=T, y=v] {figures/tuning/data/limits_caps.dat};
  \addlegendentry{$v_{\text{cap}} = D / (T(k/2+1))$}
  \addplot[sB, very thick, dashed] table[x=T, y=wz] {figures/tuning/data/limits_caps.dat};
  \addlegendentry{$\omega_{z,\text{cap}} = \mathrm{tol}/(2T)$}
  \draw[sD, thick, densely dotted] (axis cs:0.2,0.4) -- (axis cs:1.0,0.4)
    node[pos=1, font=\scriptsize, sD!30!black, anchor=west] {$v_x$ 0.4};
  \draw[sD, thick, densely dotted] (axis cs:0.2,0.3) -- (axis cs:1.0,0.3)
    node[pos=1, font=\scriptsize, sD!30!black, anchor=west] {$\omega_z$ 0.3};
  \node[font=\scriptsize, fill=white, inner sep=1pt, anchor=north east] at (axis cs:0.49,0.79) {surge $T{=}0.5$: $v_x$ 0.8};
  \node[font=\scriptsize, fill=white, inner sep=1pt, anchor=south west] at (axis cs:0.6,0.667) {sway $T{=}0.6$: $v_y$ 0.67};
  \node[font=\scriptsize, fill=white, inner sep=1pt, anchor=south west] at (axis cs:0.305,0.43) {yaw $T{=}0.3$: $\omega_z$ 0.42};
  \addplot[only marks, mark=*, sA, mark size=2.2pt] coordinates {(0.5,0.8) (0.6,0.666667)};
  \addplot[only marks, mark=square*, sB, mark size=2.2pt] coordinates {(0.3,0.416667)};
\end{groupplot}
\end{tikzpicture}
\caption{Limits from rules. (a) After a stop command the boat keeps its speed for about
one lag $T$, then the planner's deceleration ramp of $k$ lags brings it to rest; the
distance covered, $vT + v^2/2a = vT(k/2+1)$, must fit inside the stop distance $D$.
(b) The resulting caps against $T$, with a turn's overshoot $\omega_z T$ kept to half the
yaw tolerance. The marks are the rule outputs for the lags of \code{nav2\_rules}
(\code{config/characterize.yaml}): surge $T = 0.35 + 0.15$\,s gives $v_x = 0.8$\,m/s, sway
$0.45 + 0.15$\,s gives $v_y = 0.67$\,m/s, yaw $0.15 + 0.15$\,s gives $\omega_z = 0.42$\,rad/s
(\code{analysis/characterization.rule\_limits}). Each Nav2 limit is the smaller of the rule
and the measured capability. Dotted grey: today's configured $v_x = 0.4$\,m/s and
$\omega_z = 0.3$\,rad/s. Data: \code{figures/tuning/data/limits\_caps.dat}.}
\label{fig:limits}
\end{figure}


\subsection{Stopped velocities}
The precise goal checker calls the boat stopped below its stopped velocities,
\code{trans\_stopped\_velocity} and \code{rot\_stopped\_velocity}. Afloat the measured velocity never reads zero, so both sit at
least twice the per-sample neutral SD of their channel: $\ge 2\sigma_f$ and $\ge 2\sigma_r$.

\subsection{Goal tolerances: fix tier, lever arm and spin drift}
\label{sec:spin}
A goal check is a position check, so its tolerance must cover what the position does with the
boat standing still: twice the GPS noise of the fix tier in use, plus the uncorrected antenna
lever arm $l$. A third term comes from the yaw.

\paragraph{Spin drift.} Nav2 steers \code{base\_link}, but the hull turns about its centre of
mass. With the CoM a distance $d$ from base\_link, a turn in place at rate $r$ moves base\_link
on a circle of radius $d$ at speed $r d$, a sway the velocity loop reads as an error and pushes
back on one lag late. Uncorrected, a 180$^\circ$ turn carries base\_link across the chord $2d$;
measured with the loop, base\_link walks $1.2$--$2.5\,d$: $0.156$\,m at 0.3\,rad/s and $0.25$\,m at
0.8\,rad/s for $d = 0.10$\,m \ev{ad hoc Python-plant runs t3/t5, not a study row}; the range
and the factor $c$ below are read from those runs and are estimates until a study matrix
covers them. The walk grows with the rate because the
correction lags (Figure~\ref{fig:spin}). It matters most in the last turn to the goal yaw: a walk
past the xy tolerance makes the goal checker fail and the boat re-approach and turn again.
Hence
\[
  \text{xy\_goal\_tolerance} \ge 2\sigma_p + l + c\,d,\qquad c = 1.6 \text{ at } 0.3\text{\,rad/s},\; 2.0 \text{ at } 0.45\text{\,rad/s},
\]
and the in-place rotate speed stays at or under 0.3\,rad/s while $d > 0.08$\,m. Precise goals
(docking) use the RTK-fixed $\sigma_p$ only. The GPS noise itself costs hold accuracy whatever the
tuning: $\sigma_p = 0.5$\,m gives 0.6\,m hold RMS and $1.5$\,m (float) 2.3\,m
\ev{disturb\_v1:gps0.5, gps1.5\_float}.

\paragraph{Turning to the goal yaw inside the tolerance.} With \code{rotate\_to\_goal\_heading}
the rotation shim turns to the goal yaw, closed loop, once inside the xy tolerance, instead of
leaving the final yaw to MPPI's critics. In crane\_ml pair runs it cut goal overshoot from 0.24 to
0.02\,m and yaw reversals from 5.7 to 3.3\,/min for 2.3\,s of goal time
\ev{tune\_A0\_pair, tune\_A1\_pair, seeds 0--2}; it is adopted.

\begin{figure}[tbp]
\centering
\begin{tikzpicture}[>=Stealth, font=\footnotesize, x=1.12cm, y=1.12cm]
  % ---- (a) an in-place turn about the centre of mass -----------------------------------
  \def\d{1.15}
  \begin{scope}
    \node[font=\small] at (0,2.0) {(a) a $180^\circ$ turn about $C$};
    % hull before (solid) and after (dashed): base_link B at the hull centre, C a distance d aft
    \draw[thick, fill=sD!12] (\d-1.45,-0.5) -- (\d+1.0,-0.5) -- (\d+1.45,0) -- (\d+1.0,0.5)
      -- (\d-1.45,0.5) -- cycle;
    \draw[thick, dashed, sA] (-\d+1.45,-0.5) -- (-\d-1.0,-0.5) -- (-\d-1.45,0) -- (-\d-1.0,0.5)
      -- (-\d+1.45,0.5) -- cycle;
    \draw[->, thick, sA] (\d,0) arc[start angle=0, end angle=180, radius=\d];
    \fill (0,0) circle (1.8pt) node[above=1pt] {$C$};
    \fill[sB] (\d,0) circle (1.8pt) node[right=1pt, sB] {$B_0$};
    \fill[sB] (-\d,0) circle (1.8pt) node[left=1pt, sB] {$B_{180}$};
    \draw[<->, sD!30!black] (0,-0.2) -- node[below, font=\scriptsize] {$d$} (\d,-0.2);
    \draw[<->, sB, thick] (-\d,-0.8) -- node[below, font=\scriptsize] {chord $2d$} (\d,-0.8);
    \node[sA, font=\scriptsize, anchor=south] at (0,\d) {$B$ on a circle of radius $d$};
    \node[font=\scriptsize, align=center, anchor=north] at (0,-1.3)
      {lag and sway coupling: the measured walk\\ of $B$ is about $1.2$--$2.5\,d$};
  \end{scope}
  % ---- (b) the xy tolerance budget ----------------------------------------------------
  \begin{scope}[xshift=4.0cm]
    \node[font=\small, anchor=west] at (0,2.0) {(b) the goal's $xy$ tolerance must cover};
    \draw[fill=sC!25] (0,0.3) rectangle (2.4,1.0);
    \draw[fill=sB!25] (2.4,0.3) rectangle (4.0,1.0);
    \draw[fill=sA!25] (4.0,0.3) rectangle (6.6,1.0);
    \node at (1.2,0.65) {$2\sigma_p$};
    \node at (3.2,0.65) {$l$};
    \node at (5.3,0.65) {$c\,d$};
    \node[font=\scriptsize, align=center, anchor=north] at (1.2,0.2) {position noise\\ (both ends)};
    \node[font=\scriptsize, align=center, anchor=north] at (5.3,0.2) {rotation-centre walk,\\ $c \approx 1.2$--$2.5$};
    \draw[decorate, decoration={brace, amplitude=5pt}, thick] (0,1.1) -- (6.6,1.1)
      node[midway, above=5pt] {$2\sigma_p + l + c\,d \ \le\ $ \code{xy\_goal\_tolerance}};
  \end{scope}
\end{tikzpicture}
\caption{The rotation-centre offset. (a) The hull turns about its centre of mass $C$, but
the controller holds \code{base\_link} $B$, a distance $d$ away, so an in-place turn of
$180^\circ$ carries $B$ along a circle of radius $d$, a chord of $2d$. The velocity lag and
the sway the turn couples in make the walk actually seen about $1.2$--$2.5\,d$. For
$d = 0.10$\,m an ad hoc Python-plant run (\code{t3}/\code{t5}, not a study row) walked
$B$ 0.156\,m at $\omega_z = 0.3$\,rad/s and 0.25\,m at 0.8\,rad/s; the frozen study rows
(\code{com\_aft10}, \code{com\_port5}, seed~0) record only the turn's skid. (b) The $xy$ goal
tolerance must hold the position noise at both ends, $2\sigma_p$, the term $l$ and the
walk $c\,d$. Proportions illustrative.}
\label{fig:spin}
\end{figure}

\file{mhseals\_testing analysis/characterization.py (rule\_limits, nav2\_profile), config/characterize.yaml (nav2\_rules)}
